% Scope: Derive diffusion encoding from stochastic motion and gradient waveforms; progress from ADC to tensor and limited non-Gaussian models with clear identifiability boundaries.
% Status: architecture only; no chapter prose or completed problems yet.
% Prerequisites: Chapters 1, 5, 7 and 9.
% Planned worked examples: Stejskal--Tanner b-value, general b-matrix and tensor metrics.
% Planned figures: Brownian propagators, diffusion gradients and tensor ellipsoids.
% Planned challenge problems: Design equal-b unequal-time experiments; diagnose a biased tensor fit.
\chapter{Diffusion MRI}
\label{ch:10}

\section{Stochastic motion and transport}
\label{sec:10:stochastic-motion-and-transport}
\subsection{Brownian motion and the diffusion equation}
\subsection{Gaussian propagator and mean-square displacement}
\subsection{Characteristic functions and displacement encoding}

\section{Gradient-waveform encoding}
\label{sec:10:gradient-waveform-encoding}
\subsection{Stejskal--Tanner experiment and finite pulses}
\subsection{General waveform b-value and b-matrix}
\subsection{Imaging-gradient cross-terms and diffusion time}

\section{Apparent diffusion and tensor imaging}
\label{sec:10:apparent-diffusion-and-tensor-imaging}
\subsection{ADC and log-signal fitting}
\subsection{Diffusion tensor, eigenvectors and eigenvalues}
\subsection{Mean diffusivity, fractional anisotropy and invariants}

\section{Beyond Gaussian diffusion}
\label{sec:10:beyond-gaussian-diffusion}
\subsection{Restricted diffusion and time dependence}
\subsection{Diffusion kurtosis and cumulant limits}
\subsection{Multi-compartment models and IVIM}

\section{Acquisition and correction}
\label{sec:10:acquisition-and-correction}
\subsection{Spin-echo EPI and diffusion preparation}
\subsection{Motion, eddy currents and susceptibility distortion}
\subsection{Gradient nonlinearity, b-vector rotation and noise-floor bias}

\section{Interpretation and advanced encoding}
\label{sec:10:interpretation-and-advanced-encoding}
\subsection{Acute ischemia and oncologic examples}
\subsection{Tractography, crossing fibers and inference limits}
\subsection{Oscillating gradients and tensor-valued encoding}

\section{Worked examples}
\label{sec:10:worked-examples}
% Planned calculations: Stejskal--Tanner b-value, general b-matrix and tensor metrics.
\subsection{Derivation and quantitative calculation}
\subsection{Assumptions, limiting cases and scanner interpretation}

\section{Challenge problems}
\label{sec:10:challenge-problems}
% Problem design: Design equal-b unequal-time experiments; diagnose a biased tensor fit.
% Use challengeproblem and stable labels prob:10:<topic>.
% Complete solutions belong in Appendix E, section for this chapter.
% Use \begin{solution}{prob:10:<topic>} ... \end{solution} there.
\subsection{Derivation and quantitative design}
\subsection{Model criticism and integrated reasoning}
